\section{Proof of Theorem A}
\label{sec:main-theorem}

\subsection{The first six-generated family}
\label{sec:family-one}

In this subsection, we treat the first family \eqref{eq:family-one} in the Introduction. Put
$$
\begin{aligned}
 I={}&( g_1 = y^2-x^{b-d}z,\ g_2 = x^{2b+2-e},\ g_3 = x^{b+1}y,\
        g_4 = x^{b-d+1}z,\ g_5 = yz,\ g_6 = z^2),\\
 &\hspace{35mm} b>d\geq0,\qquad e\in\{0,1\}.
\end{aligned}
$$

\subsubsection{Cohen--Macaulayness}

\begin{proposition}
\label{prop:family-one-local-reductions}
Define $Q\subseteq I$ by
\[
\begin{array}{c|c|l}
 e&d&Q\\ \hline
 1&0,1&(g_1,g_5,g_2+g_6)\\
 1&d\geq2&(g_1,g_2+g_5,g_4+g_6)\\
 0&0&(g_1+g_5,g_3+g_5,g_2+g_4+g_5+g_6)\\
 0&d\geq1&(g_1,g_3+g_5,g_2+g_4+g_6).
\end{array}
\]
Then
\[
 I^2=QI+\m I^2.
\]
Consequently, $\calR(I)$ is Cohen--Macaulay by Proposition~\ref{prop:local-reduction-criterion}.
\end{proposition}

\begin{proof}
	We consider separately the four cases in the definition of $Q$.
	
	\medskip
	
	\noindent
	\underline{$e=1$, $d =0, 1$}: Since $Q=(g_1, g_5, g_2+g_6)$, we have $I=Q+ (g_3, g_4, g_6)$.
	Hence $I^2=QI + (g_ig_j \mid i,j \in \{3,4,6\})$.
	Since $d\in\{0,1\}$, all the powers of $x$ occurring below are non-negative.
	\begin{itemize}
		\item $g_4g_6=-xg_1g_6+xg_5^2\in QI+\m I^2$.
		\item $g_3^2=xg_1g_2+g_4(g_2+g_6)-g_4g_6\in QI+\m I^2$.
		\item $
		 g_3g_4=x^{1-d}g_5(g_2+g_6)-x^{1-d}g_5g_6\in QI.
		$
		\item $
		 g_3g_6=x^dg_4g_5\in QI.
		$
		\item $
		 g_4^2=-xg_1g_4+x^{1-d}g_3g_5\in QI.
		$
		\item $
		 g_2g_6=-x^{2d}g_1g_4+x^dg_3g_5\in QI,
		$
		we obtain
		$
		 g_6^2=(g_2+g_6)g_6-g_2g_6\in QI.
		$
	\end{itemize}
	Thus all the six products belong to $QI+\m I^2$, and hence
	$I^2=QI+\m I^2$.

	\medskip
	
	\noindent
	\underline{$e=1$, $d\geq2$}: Since $Q=(g_1,g_2+g_5,g_4+g_6)$, we have $I=Q+(g_3,g_5,g_6)$.
	Hence $I^2=QI+(g_ig_j\mid i,j\in\{3,5,6\})$.
	Since $d\geq2$, all the powers of $x$ occurring below are positive.
	\begin{itemize}
		\item $g_3^2=xg_1g_2+(g_4+g_6)g_2-x^{2d-1}g_4^2\in QI+\m I^2$.
		\item $g_3g_5=x^dg_1g_4+x^{d-1}g_4^2\in QI+\m I^2$.
		\item $g_3g_6=x^dg_4g_5\in\m I^2$.
		\item $g_5^2=(g_2+g_5)g_5-x^{d-1}g_3g_4\in QI+\m I^2$.
		\item $g_5g_6=(g_2+g_5)g_6-x^{2d-1}g_4^2\in QI+\m I^2$.
		\item $g_6^2=(g_4+g_6)g_6-x(g_5^2-g_1g_6)\in QI+\m I^2$.
	\end{itemize}
	Thus all the six products belong to $QI+\m I^2$, and hence
	$I^2=QI+\m I^2$.

	\medskip
	
	\noindent
	\underline{$e=0$, $d=0$}: Since
	$Q=(g_1+g_5,g_3+g_5,g_2+g_4+g_5+g_6)$, we have
	$I=Q+(g_2,g_4,g_5)$.
	Hence $I^2=QI+(g_ig_j\mid i,j\in\{2,4,5\})$.
	We first note that
	$g_2g_6=-xg_1g_4+xg_3g_5\in\m I^2$ and
	$g_4g_6=-xg_1g_6+xg_5^2\in\m I^2$.
	\begin{itemize}
		\item $g_2g_4=x(g_3^2-g_1g_2)\in\m I^2$.
		\item $g_4^2=g_2g_6\in\m I^2$.
		\item $g_4g_5=(g_2+g_4+g_5+g_6)g_4-g_2g_4-g_4^2-g_4g_6\in QI+\m I^2$.
		\item Since $g_3g_4=g_2g_5$, we have
		$g_2g_5=(g_3+g_5)g_4-g_4g_5\in QI+\m I^2$.
		\item Since $g_3g_6=g_4g_5$, we have
		$g_5^2=(g_2+g_4+g_5+g_6)g_5-g_2g_5-(g_3+g_5)g_6\in QI+\m I^2$.
		\item $g_2^2=(g_2+g_4+g_5+g_6)g_2-g_2g_4-g_2g_5-g_2g_6\in QI+\m I^2$.
	\end{itemize}
	Thus all the six products belong to $QI+\m I^2$, and hence
	$I^2=QI+\m I^2$.

	\medskip
	
	\noindent
	\underline{$e=0$, $d\geq1$}: Since $Q=(g_1,g_3+g_5,g_2+g_4+g_6)$, we have
	$I=Q+(g_2,g_4,g_5)$.
	Hence $I^2=QI+(g_ig_j\mid i,j\in\{2,4,5\})$.
	We first note that
	$g_2g_6=-x^{2d+1}g_1g_4+x^{d+1}g_3g_5\in\m I^2$ and
	$g_4g_6=-xg_1g_6+xg_5^2\in\m I^2$.
	\begin{itemize}
		\item $g_2g_4=x(g_3^2-g_1g_2)\in\m I^2$.
		\item $g_2g_5=x^dg_3g_4\in\m I^2$.
		\item $g_4^2=(g_2+g_4+g_6)g_4-g_2g_4-g_4g_6\in QI+\m I^2$.
		\item Since $g_3g_6=x^dg_4g_5\in\m I^2$, we have
		$g_5g_6=(g_3+g_5)g_6-g_3g_6\in QI+\m I^2$, and hence
		$g_4g_5=(g_2+g_4+g_6)g_5-g_2g_5-g_5g_6\in QI+\m I^2$.
		\item $g_5^2=(g_3+g_5)g_5-x^dg_1g_4-x^{d-1}g_4^2\in QI+\m I^2$.
		\item $g_2^2=(g_2+g_4+g_6)g_2-g_2g_4-g_2g_6\in QI+\m I^2$.
	\end{itemize}
	Thus all the six products belong to $QI+\m I^2$, and hence
	$I^2=QI+\m I^2$.
\end{proof}

\begin{remark}
Proposition~\ref{prop:family-one-local-reductions} shows that $QS_\m$ is a
reduction of $IS_\m$; it does not assert that $Q$ itself is a reduction of
$I$ in $S$.  Since no global reduction is needed for the proof of
Theorem~\ref{thm:main}, we do not pursue one here.
\end{remark}

\subsubsection{Normality}

We next prove that $I$ is normal.
Recall from Subsection~\ref{subsec:multiplicity-three} that
$\deg x=3$, $\deg y=3b+2-e$, and $\deg z=3b+3d+4-2e$; in particular,
the least relation degree is $\ell=6b+4-2e$.

Let $f=y^2$ and $J=I+(f)$. Since $\varphi_H(f) = t^{6b+4-2e} = t^\ell$, 
$$
J=(y^2, x^{b-d}z, x^{2b+2-e}, 	x^{b+1}y, yz, z^2)
$$
is a monomial ideal with six minimal generators.
As observed in Section~\ref{sec:target-ideals}, $J=I_\ell$ and hence $J$ is integrally closed.  Since $J$ is a monomial ideal with six minimal generators, it is normal by \cite{AtakaMatsuoka2026}.

\begin{lemma}\label{lem:common_firstfamily}
The following assertions hold.
\begin{enumerate}
	\item $\m f \subseteq I$.
	\item $g_6 f \in I^2$.
\end{enumerate}
	
\end{lemma}

\begin{proof}
	(1) $y^2z = y(yz) \in I$ is obvious.
	Furthermore, $xy^2 = xg_1 + g_4$ and $y^3 = yg_1 + x^{b-d}(yz)$ imply $xy^2, y^3 \in I$.
	
	(2) It is easy that $g_6 f = y^2z^2 = (yz)^2 \in I^2$.
\end{proof}

We now apply the strategy described in Subsection~\ref{subsec:normality}.

Suppose that $I^s$ is not integrally closed for some positive integer $s$.
Since $\ol{I^s}$ is homogeneous, take a homogeneous element
$\xi\in\ol{I^s}\setminus I^s$ of degree $\delta$.  Since
$\xi\in\ol{I^s}\subseteq\ol{J^s}=J^s$, we may choose the least positive
integer $m$ such that
$$
\xi \in I^s + fI^{s-1} + \cdots + f^m I^{s-m}
$$
and write
$$
\xi = \xi_0 + f^m \omega
$$
with $\xi_0\in I^s+fI^{s-1}+\cdots+f^{m-1}I^{s-m+1}$ and
$\omega\in I^{s-m}$.  We may assume that $\omega$ is homogeneous of degree
$\delta-m\ell$.  By Lemma~\ref{lem:common_firstfamily}(2) and the reductions
in Subsection~\ref{subsec:normality}, we may put
$$
\Lambda_0=\{\lambda\in\Lambda^{s-m}_{\delta-m\ell}\mid\lambda_6=0\}
$$
and write
$$
\omega = \sum_{\lambda \in \Lambda_0} c_\lambda g^\lambda
$$
with $c_\lambda\in k$.


\begin{lemma}\label{lem:firstfamily_cases}
	The following assertions hold.
	\begin{enumerate}
		\item If $e=0$, then $fg_2 \in I^2$.
		\item Suppose $d=0$. Then $fg_4 \in I^2$. Furthermore, if $e=1$, then $fg_3 \in I^2$.
		\item Suppose $d=1$. Then $g_3g_5 \equiv g_4^2 \pmod{\m I^2}$. Furthermore, if $e=1$, then $g_3^2 \equiv g_2g_4 \pmod{\m I^2}$, $g_2g_5 = g_3g_4$, and
		$g_3^3\equiv g_2^2g_5\pmod{\m I^3}$.
		\item Suppose $d\geq2$. Then $g_3g_5\in\m I^2$.  If $e=1$, then $g_2g_5\in\m I^2$ and $g_3^2\equiv g_2g_4\pmod{\m I^2}$.
	\end{enumerate}
\end{lemma}

\begin{proof}
The assertions follow from the following identities:
\begin{align*}
 fg_2&=g_3^2 &&(e=0),\\
 fg_4&=g_3g_5 &&(d=0),\\
 fg_3&=g_1g_3+g_2g_5 &&(d=0,\ e=1),\\
 g_4^2-g_3g_5&=-xg_1g_4 &&(d=1),\\
 g_3^2-g_2g_4&=xg_1g_2 &&(d=1,\ e=1),\\
 g_2g_5&=g_3g_4 &&(d=1,\ e=1),\\
 g_3^3-g_2^2g_5&=xg_1g_2g_3 &&(d=1,\ e=1),\\
 g_3g_5&=x^dg_1g_4+x^{d-1}g_4^2 &&(d\geq2),\\
 g_2g_5&=x^{d-1}g_3g_4 &&(d\geq2,\ e=1),\\
 g_3^2-g_2g_4&=xg_1g_2 &&(d\geq2,\ e=1).
\end{align*}
\end{proof}

Using Lemma~\ref{lem:firstfamily_cases}, we can impose the
following additional conditions on the elements of $\Lambda_0$.  In each
case, let $\Lambda$ denote the resulting set.  We may then write
$$\omega=\sum_{\lambda\in\Lambda}c_\lambda g^\lambda.$$

For $d=0$, the products involving $f$ in Lemma~\ref{lem:firstfamily_cases}
give the first two rows below.  Suppose that $d=1$ and $e=1$.  We first use
$g_4^2\equiv g_3g_5$ to arrange $0\leq\lambda_4\leq1$.  If $\lambda_4=1$
and either $\lambda_2>0$ or $\lambda_3>0$, the other two quadratic relations remove
$g_4$.  When $\lambda_4=0$, the cubic relation allows us to arrange
$0\leq\lambda_3\leq2$.  If $d=1$ and $e=0$, the equality $fg_2=g_3^2$
removes $g_2$, and $g_4^2\equiv g_3g_5$ again gives
$0\leq\lambda_4\leq1$.  Finally, suppose that $d\geq2$.  If $e=1$, use
$g_3^2\equiv g_2g_4$ to arrange $0\leq\lambda_3\leq1$; when $\lambda_5>0$,
the two products belonging to $\m I^2$ remove $g_2$ and $g_3$.  If $e=0$,
the equality $fg_2=g_3^2$ removes $g_2$, and $g_3g_5\in\m I^2$ removes
$g_3$ whenever $\lambda_5>0$.

\[
\begin{array}{c|l}
 \text{Case}&\text{Additional Conditions} \\ \hline \hline
 d=0,\, e=1 & \lambda_3 = \lambda_4 = 0\\
 \hline
 d=0,\, e=0 & \lambda_2 = \lambda_4 = 0\\
 \hline
 d=1,\, e=1 & \begin{gathered}0\leq\lambda_3\leq2,\quad 0\leq\lambda_4\leq1,\\
                    \text{if }\lambda_4=1,\text{ then }\lambda_2=\lambda_3=0
                    \end{gathered}\\
 \hline
 d=1,\, e=0 & \lambda_2 = 0, \, 0 \le \lambda_4 \le 1\\
 \hline
 d\ge 2, \, e=1 & \begin{gathered}0\leq\lambda_3\leq1,\\
                    \text{if }\lambda_5>0,\text{ then }\lambda_2=\lambda_3=0
                    \end{gathered}\\
 \hline
 d\ge 2, \, e=0 & \begin{gathered}\lambda_2=0,\\
                    \text{if }\lambda_5>0,\text{ then }\lambda_3=0
                    \end{gathered}
\end{array}
\]


We now prove, in each of the six cases, that
$c_\lambda=0$ for every $\lambda\in\Lambda$.  For a homomorphism
$\psi_\varepsilon\colon S\to k[[t]]$, we write
$\ord_\varepsilon(g)=\ord_t\psi_\varepsilon(g)$ for $g\in S$.


\medskip

\noindent
\underline{$d=0$, $e=1$}: 
For each $\varepsilon \in k^{\times}$, let $\psi_\varepsilon : S \to k[[t]]$ be a ring homomorphism defined by
$$
x \mapsto t^3,\quad y \mapsto \varepsilon t^{3b+1}, \quad z \mapsto \varepsilon^2 t^{3b+2} - t^{3b+3}.
$$
We note the following:
\[
\begin{array}{c|c|c}
 \text{Element}&\psi_\varepsilon&\ord_\varepsilon\\ \hline
 f=y^2&\varepsilon^2t^{6b+2}&6b+2\\
 g_1=y^2-x^bz&t^{6b+3}&6b+3\\
 g_2=x^{2b+1}&t^{6b+3}&6b+3\\
 g_5=yz&\varepsilon^3t^{6b+3}-\varepsilon t^{6b+4}&6b+3
\end{array}
\]

Put $\nu=6b+3$. The omitted generators $g_3,g_4,g_6$ have
orders $\nu+1,\nu+2,\nu+1$, respectively. Hence
$\psi_\varepsilon(I)k[[t]]\subseteq(t^\nu)$. Set
\[
 G=\{i\mid\ord_t\psi_\varepsilon(g_i)=\nu\}=\{1,2,5\}.
\]
Every $\lambda\in\Lambda$ has support in $G$.
By Lemma~\ref{lem:coefficient-separation}, it suffices to check that
$\lambda\mapsto3\lambda_5$ is injective on $\Lambda$.

For a fixed integer $i$, an element
$\lambda\in\Lambda$ with $\lambda_5=i$ is determined by the following two
equations in $\lambda_1$ and $\lambda_2$:
$$
\begin{cases}
	\lambda_1 + \lambda_2 = s-m-i, \\
	\lambda_1 \deg g_1 + \lambda_2 \deg g_2
	= \delta-m\ell-i\deg g_5.
\end{cases}
$$
Since $\deg g_2-\deg g_1=1$, this system has at most one
solution. Thus the required map is injective and
Lemma~\ref{lem:coefficient-separation} gives $c_\lambda=0$ for every
$\lambda\in\Lambda$.


\medskip

\noindent
\underline{$d=0$, $e=0$}: 
For each $\varepsilon\in k^\times$, define
$\psi_\varepsilon\colon S\to k[[t]]$ by
$$
x\longmapsto t^2,\qquad
y\longmapsto t^{2b+1},\qquad
z\longmapsto t^{2b+2}-\varepsilon t^{2b+3}.
$$
Put $\nu=4b+3$.  The elements that occur in $\omega$ have the following
images:
\[
\begin{array}{c|c|c}
 \text{Element}&\psi_\varepsilon&\ord_\varepsilon\\ \hline
 f=y^2&t^{4b+2}&\nu-1\\
 g_1=y^2-x^bz&\varepsilon t^{4b+3}&\nu\\
 g_3=x^{b+1}y&t^{4b+3}&\nu\\
 g_5=yz&t^{4b+3}-\varepsilon t^{4b+4}&\nu
\end{array}
\]
The omitted generators $g_2,g_4,g_6$ all have order $\nu+1$.
Thus $\psi_\varepsilon(I)k[[t]]\subseteq(t^\nu)$. Set
\[
 G=\{i\mid\ord_t\psi_\varepsilon(g_i)=\nu\}=\{1,3,5\}.
\]
Every $\lambda\in\Lambda$ has support in $G$. Lemma~\ref{lem:coefficient-separation} reduces the claim to
injectivity of $\lambda\mapsto\lambda_1$ on $\Lambda$.

Fix $i$.  An element $\lambda\in\Lambda$ with $\lambda_1=i$ is determined
by
$$
\begin{cases}
 \lambda_3+\lambda_5=s-m-i,\\
 \lambda_3\deg g_3+\lambda_5\deg g_5
 =\delta-m\ell-i\deg g_1.
\end{cases}
$$
Since $\deg g_5-\deg g_3=1$, this system has at most one solution.
Therefore all the coefficients $c_\lambda$ vanish.

\medskip

\noindent
\underline{$d=1$, $e=1$}:
For $\varepsilon\in k^\times$, let $\psi_\varepsilon\colon S\to k[[t]]$
be given by
$$
x\longmapsto t,\qquad
y\longmapsto t^b,\qquad
z\longmapsto t^{b+1}-\varepsilon t^{b+2}.
$$
With $\nu=2b+1$, we have
\[
\begin{array}{c|c|c}
 \text{Element}&\psi_\varepsilon&\ord_\varepsilon\\ \hline
 f=y^2&t^{\nu-1}&\nu-1\\
 g_1&\varepsilon t^\nu&\nu\\
 g_2&t^\nu&\nu\\
 g_3&t^\nu&\nu\\
 g_4&t^\nu-\varepsilon t^{\nu+1}&\nu\\
 g_5&t^\nu-\varepsilon t^{\nu+1}&\nu
\end{array}
\]
The remaining generator $g_6$ has order $\nu+1$, so
$\psi_\varepsilon(I)k[[t]]\subseteq(t^\nu)$. Set
\[
 G=\{i\mid\ord_t\psi_\varepsilon(g_i)=\nu\}=\{1,2,3,4,5\}.
\]
Every $\lambda\in\Lambda$ has support in $G$.
We check the injectivity of $\lambda\mapsto\lambda_1$ on $\Lambda$
to apply Lemma~\ref{lem:coefficient-separation}.


First suppose that $\lambda_4=0$.  After fixing $\lambda_1=i$, the remaining
exponents satisfy
$$
\lambda_2+\lambda_3+\lambda_5=s-m-i
$$
and
$$
\lambda_3+3\lambda_5
=\delta-m\ell-i\deg g_1-(6b+3)(s-m-i).
$$
The right-hand side is fixed.  Since $0\leq\lambda_3\leq2$, its residue
modulo $3$ determines $\lambda_3$, after which $\lambda_5$ and $\lambda_2$
are determined in this order.  Thus there is at most one such element of
$\Lambda$.

If $\lambda_4=1$, then $\lambda_2=\lambda_3=0$, and the number of factors
determines $\lambda_5$.  It remains to see that the two forms cannot occur
with the same value of $i$ and the same weighted degree.  Suppose that
$(i,\lambda_2,\lambda_3,0,\lambda_5,0)$ and
$(i,0,0,1,\mu_5,0)$ did so.  The number of factors gives
$\mu_5=\lambda_2+\lambda_3+\lambda_5-1$.  Comparing their weighted degrees
then yields
$$
3\lambda_2+2\lambda_3=1,
$$
which has no solution in non-negative integers. Hence the required
map is injective, and Lemma~\ref{lem:coefficient-separation} gives
$c_\lambda=0$ for every $\lambda\in\Lambda$.


\medskip

\noindent
\underline{$d=1$, $e=0$}:
Use the same map as in the preceding case, again with $\nu=2b+1$.  For the
elements occurring in $\omega$, we have
\[
\begin{array}{c|c|c}
 \text{Element}&\psi_\varepsilon&\ord_\varepsilon\\ \hline
 f=y^2&t^{\nu-1}&\nu-1\\
 g_1&\varepsilon t^\nu&\nu\\
 g_3&t^\nu&\nu\\
 g_4&t^\nu-\varepsilon t^{\nu+1}&\nu\\
 g_5&t^\nu-\varepsilon t^{\nu+1}&\nu
\end{array}
\]
The omitted generators $g_2,g_6$ have order $\nu+1$. Hence
$\psi_\varepsilon(I)k[[t]]\subseteq(t^\nu)$. Set
\[
 G=\{i\mid\ord_t\psi_\varepsilon(g_i)=\nu\}=\{1,3,4,5\}.
\]
Every $\lambda\in\Lambda$ has support in $G$.
We apply Lemma~\ref{lem:coefficient-separation} by checking
injectivity of $\lambda\mapsto\lambda_1$ on $\Lambda$.

Fix $\lambda_1=i$.  Since $\lambda_2=\lambda_6=0$, the number of factors is
$$
\lambda_3+\lambda_4+\lambda_5=s-m-i.
$$
The weighted-degree equation, after subtracting
$(6b+5)(s-m-i)$ and dividing by $2$, has the form
$$
\lambda_4+2\lambda_5=c
$$
for a fixed integer $c$.  Since $0\leq\lambda_4\leq1$, the parity of $c$
determines $\lambda_4$, and then $\lambda_5$ and $\lambda_3$ are uniquely
determined.  Thus every coefficient $c_\lambda$ is zero.

\medskip

\noindent
\underline{$d\geq2$, $e=1$}:
We first use the homomorphism $\psi_\varepsilon\colon S\to k[[t]]$ defined by
$$
x\longmapsto t,\qquad
y\longmapsto t^b,\qquad
z\longmapsto t^{b+d}-\varepsilon t^{b+d+1}.
$$
Put $\nu_1=2b+1$.  Then
\[
\begin{array}{c|c|c}
 \text{Element}&\psi_\varepsilon&\ord_\varepsilon\\ \hline
 f=y^2&t^{\nu_1-1}&\nu_1-1\\
 g_1&\varepsilon t^{\nu_1}&\nu_1\\
 g_2&t^{\nu_1}&\nu_1\\
 g_3&t^{\nu_1}&\nu_1\\
 g_4&t^{\nu_1}-\varepsilon t^{\nu_1+1}&\nu_1\\
 g_5&t^{\nu_1+d-1}-\varepsilon t^{\nu_1+d}&\nu_1+d-1
\end{array}
\]
The remaining generator $g_6$ has order $2b+2d>\nu_1$.
Thus $\psi_\varepsilon(I)k[[t]]\subseteq(t^{\nu_1})$. Set
\[
 G_1=\{i\mid\ord_t\psi_\varepsilon(g_i)=\nu_1\}=\{1,2,3,4\}.
\]
The elements of $\Lambda$ with support in $G_1$ are precisely those with $\lambda_5=0$. For this subset,
Lemma~\ref{lem:coefficient-separation} requires injectivity of
$\lambda\mapsto\lambda_1$.

After fixing $\lambda_1=i$, the remaining exponents satisfy
$$
\lambda_2+\lambda_3+\lambda_4=s-m-i
$$
and
$$
\lambda_3+2\lambda_4
=\delta-m\ell-i\deg g_1-(6b+3)(s-m-i).
$$
Since $0\leq\lambda_3\leq1$, the parity of the right-hand side determines
$\lambda_3$, and the two equations then determine $\lambda_4$ and
$\lambda_2$.  Hence the coefficients with $\lambda_5=0$ vanish.

It remains to treat the elements of $\Lambda$ with $\lambda_5>0$.  For these
elements, $\lambda_2=\lambda_3=0$.  Define a second homomorphism by
$$
x\longmapsto t,\qquad
y\longmapsto t^{b-d+1},\qquad
z\longmapsto t^{b-d+2}-\varepsilon t^{b-d+3}.
$$
For $\nu_2=2b-2d+3$, we have
\[
\begin{array}{c|c|c}
 \text{Element}&\psi_\varepsilon&\ord_\varepsilon\\ \hline
 f=y^2&t^{\nu_2-1}&\nu_2-1\\
 g_1&\varepsilon t^{\nu_2}&\nu_2\\
 g_4&t^{\nu_2}-\varepsilon t^{\nu_2+1}&\nu_2\\
 g_5&t^{\nu_2}-\varepsilon t^{\nu_2+1}&\nu_2
\end{array}
\]
For the omitted generators, the orders exceed $\nu_2$ by
\[
 \ord_\varepsilon(g_2)-\nu_2=2d-2,\qquad
 \ord_\varepsilon(g_3)-\nu_2=d-1,\qquad
 \ord_\varepsilon(g_6)-\nu_2=1.
\]
Hence $\psi_\varepsilon(I)k[[t]]\subseteq(t^{
\nu_2})$. For this second map, set
\[
 G_2=\{i\mid\ord_t\psi_\varepsilon(g_i)=\nu_2\}=\{1,4,5\}.
\]
 Every remaining exponent vector has support in
$G_2$, so the two subsets cover $\Lambda$. Apply
Lemma~\ref{lem:coefficient-separation} to the remaining sum.

For a fixed value $\lambda_1=i$, the exponents $\lambda_4$ and $\lambda_5$
satisfy
$$
\begin{cases}
 \lambda_4+\lambda_5=s-m-i,\\
 \lambda_4\deg g_4+\lambda_5\deg g_5
 =\delta-m\ell-i\deg g_1.
\end{cases}
$$
Since $\deg g_5-\deg g_4=3d-2\neq0$, the solution is unique if it exists.
The remaining coefficients therefore vanish.

\medskip

\noindent
\underline{$d\geq2$, $e=0$}:
Use first the map from the preceding case with $\nu_1=2b+1$.
Here $g_2$ has order $\nu_1+1$, and $g_6$ has order
$2b+2d>\nu_1$. Thus $\psi_\varepsilon(I)k[[t]]\subseteq(t^{\nu_1})$. Set
\[
 G_1=\{i\mid\ord_t\psi_\varepsilon(g_i)=\nu_1\}=\{1,3,4\}.
\]
The elements of $\Lambda$ with support in $G_1$ are precisely those with $\lambda_5=0$. We check injectivity of $\lambda\mapsto\lambda_1$
on this subset for Lemma~\ref{lem:coefficient-separation}.

After fixing $\lambda_1=i$, the remaining exponents satisfy
$$
\begin{cases}
 \lambda_3+\lambda_4=s-m-i,\\
 \lambda_3\deg g_3+\lambda_4\deg g_4
 =\delta-m\ell-i\deg g_1.
\end{cases}
$$
Here $\deg g_4-\deg g_3=2$, so there is at most one solution.  Hence all the
coefficients with $\lambda_5=0$ vanish.

Every remaining element of $\Lambda$ has $\lambda_5>0$ and therefore
$\lambda_3=0$. Apply the second map from the preceding case. Now
\[
 \ord_\varepsilon(g_2)-\nu_2=2d-1,\qquad
 \ord_\varepsilon(g_3)-\nu_2=d-1,\qquad
 \ord_\varepsilon(g_6)-\nu_2=1,
\]
so $\psi_\varepsilon(I)k[[t]]\subseteq(t^{\nu_2})$. For this second map, set
\[
 G_2=\{i\mid\ord_t\psi_\varepsilon(g_i)=\nu_2\}=\{1,4,5\}.
\]

Every remaining exponent vector has support in $G_2$, and the two
subsets therefore cover $\Lambda$. Apply Lemma~\ref{lem:coefficient-separation}
to the remaining sum. For fixed $\lambda_1=i$, the exponents
$\lambda_4$ and $\lambda_5$ satisfy the same
number-of-factors equation as above and the corresponding weighted-degree
equation.  Since $\deg g_5-\deg g_4=3d-1\neq0$, the solution is unique if it
exists.  Thus the remaining coefficients vanish as well.

In all six cases, every coefficient $c_\lambda$ is zero.  Hence $\omega=0$,
contrary to the minimality of $m$.  It follows that $I^s$ is integrally
closed for every $s\geq1$, and therefore $I$ is normal.

Combining this with Proposition~\ref{prop:family-one-local-reductions}, we obtain the following.

\begin{theorem}
\label{thm:family-one}
Let $I$ be an ideal in the family~\eqref{eq:family-one}.  Then the Rees algebra $\calR(I)$ is a Cohen--Macaulay normal domain.
\end{theorem}
